\section{宣统转型证据链事件锚点}
以下锚点补足革命、铁路、议政、谈判和退位的前后记录，保留中央时间与各地行政、财政和消息传播节奏之间的不一致。
\subsection{1900—1909}
\eventanchor{EQ-1909-PROVINCIAL-ASSEMBLIES-DOCUMENTS}{围绕「各省咨议局开办，地方精英…}宣统元年（1909年），围绕「各省咨议局开办，地方精英获得新的公开政治场域」的上谕、奏折与部议在中央—地方行政链中流转。核心取证：《宣统政纪》宣统元年相关条；宫中朱批奏折、内阁大库题本与《清史稿》相关志传。这一锚点记录的是文书链：它能证明呈报、上谕、部议或照会进入机构流程，不能由此推出实际执行。来源保留：此行仅确认相关文书链和可追查的行政动作；不得由其直接推断政策有效、地方服从、民众态度或单一因果。
\eventanchor{EQ-1909-RAILWAY-RECOVERY-MOVEMENT-DOCUMENTS}{围绕「铁路国有化和路权问题引发…}宣统元年（1909年），围绕「铁路国有化和路权问题引发各省商绅、公司和官府的争论」的经费、税收、赈济或工程呈报经由户部与地方衙门处理。核心取证：《宣统政纪》宣统元年相关条；户部奏销、盐政、河工或海关档案。这一锚点记录的是文书链：它能证明呈报、上谕、部议或照会进入机构流程，不能由此推出实际执行。来源保留：此行仅确认相关文书链和可追查的行政动作；不得由其直接推断政策有效、地方服从、民众态度或单一因果。
\eventanchor{EQ-1909-RAILWAY-RECOVERY-MOVEMENT-PRELUDE}{「铁路国有化和路权问题引发各省…}宣统元年（1909年），「铁路国有化和路权问题引发各省商绅、公司和官府的争论」之前的条件、信息和争议已通过中央与地方材料显现，构成该事件可追踪的前史节点。核心取证：《宣统政纪》宣统元年相关条；户部奏销、盐政、河工或海关档案；同时期地方档案、地方志、日记或对方材料应按具体地区补检。这一锚点界定可回查的前史条件；条件、传闻或信息累积不应被写成已经证实的单一直接原因。来源保留：官修编年和地方材料的记录密度与立场并不相同；本行不将条件、传闻、呈报或后见解释直接等同为确定因果。
\eventanchor{EQ-1909-RAILWAY-RECOVERY-MOVEMENT-DECISION}{围绕「铁路国有化和路权问题引发…}宣统元年（1909年），围绕「铁路国有化和路权问题引发各省商绅、公司和官府的争论」的命令、调度、照会或呈报形成独立的中央—地方行动文书链。核心取证：《宣统政纪》宣统元年相关条；户部奏销、盐政、河工或海关档案。这一锚点定位决策或调度动作；命令发出、地方接收、执行过程和实际结果仍须分别寻找证据。来源保留：奏折、上谕、部议、军机档或条约可证明行政动作和机构立场，不自动证明所有地区、群体或对手按其意图行动。
\eventanchor{EQ-1909-RAILWAY-RECOVERY-MOVEMENT-AFTERMATH}{「铁路国有化和路权问题引发各省…}宣统元年（1909年），「铁路国有化和路权问题引发各省商绅、公司和官府的争论」引发的善后、执行或地方回应构成可由不同材料追踪的后续节点。核心取证：《宣统政纪》宣统元年相关条；户部奏销、盐政、河工或海关档案；相关地区的档案、志书、契约、报刊或对方材料。这一锚点追踪后续影响；不同地区、群体与时段的反应不能由战后或改革后的概括性记忆代替。来源保留：战后、改革后或条约后的记忆、地方记录和官方总结常具有事后选择性；后续影响须按地点、群体和时间分层，不作整体化推断。
\subsection{1910—1919}
\eventanchor{EQ-1910-ZIZHENG-YUAN-DOCUMENTS}{围绕「资政院开院，中央层面的立…}宣统二年（1910年），围绕「资政院开院，中央层面的立宪咨询机构开始运作」的上谕、奏折与部议在中央—地方行政链中流转。核心取证：《宣统政纪》宣统二年相关条；宫中朱批奏折、内阁大库题本与《清史稿》相关志传。这一锚点记录的是文书链：它能证明呈报、上谕、部议或照会进入机构流程，不能由此推出实际执行。来源保留：此行仅确认相关文书链和可追查的行政动作；不得由其直接推断政策有效、地方服从、民众态度或单一因果。
\eventanchor{EQ-1910-FAMINE-AND-UNREST-DOCUMENTS}{围绕「灾荒、物价和财政压力持续…}宣统二年（1910年），围绕「灾荒、物价和财政压力持续影响晚清地方社会与政治动员」的地方呈报、案件或赈务记录将社会反应带入官府文书链。核心取证：《宣统政纪》宣统二年相关条；地方志、督抚奏折及地方档案。这一锚点记录的是文书链：它能证明呈报、上谕、部议或照会进入机构流程，不能由此推出实际执行。来源保留：此行仅确认相关文书链和可追查的行政动作；不得由其直接推断政策有效、地方服从、民众态度或单一因果。
\eventanchor{EQ-1911-RAILWAY-NATIONALIZATION-DOCUMENTS}{围绕「清廷宣布铁路干线国有化并…}宣统三年（1911年），围绕「清廷宣布铁路干线国有化并以外债筹资，保路运动迅速扩大」的上谕、奏折与部议在中央—地方行政链中流转。核心取证：《宣统政纪》宣统三年相关条；宫中朱批奏折、内阁大库题本与《清史稿》相关志传。这一锚点记录的是文书链：它能证明呈报、上谕、部议或照会进入机构流程，不能由此推出实际执行。来源保留：此行仅确认相关文书链和可追查的行政动作；不得由其直接推断政策有效、地方服从、民众态度或单一因果。
\eventanchor{EQ-1911-SICHUAN-RAILWAY-PROTEST-DOCUMENTS}{围绕「四川保路运动升级，赵尔丰…}宣统三年（1911年），围绕「四川保路运动升级，赵尔丰处置引发更大政治危机」的地方呈报、案件或赈务记录将社会反应带入官府文书链。核心取证：《宣统政纪》宣统三年相关条；地方志、督抚奏折及地方档案。这一锚点记录的是文书链：它能证明呈报、上谕、部议或照会进入机构流程，不能由此推出实际执行。来源保留：此行仅确认相关文书链和可追查的行政动作；不得由其直接推断政策有效、地方服从、民众态度或单一因果。
\eventanchor{EQ-1911-WUCHANG-UPRISING-DOCUMENTS}{清廷围绕「武昌起义爆发，新军与…}宣统三年（1911年），清廷围绕「武昌起义爆发，新军与革命团体控制武汉部分地区」调遣军队、筹措饷械并处理战报，中央与地方的军政文书形成连续记录。核心取证：《宣统政纪》宣统三年相关条；军机处档、战事方略及地方奏报。这一锚点记录的是文书链：它能证明呈报、上谕、部议或照会进入机构流程，不能由此推出实际执行。来源保留：此行仅确认相关文书链和可追查的行政动作；不得由其直接推断政策有效、地方服从、民众态度或单一因果。
\eventanchor{EQ-1911-PROVINCIAL-INDEPENDENCE-DOCUMENTS}{围绕「多省宣布脱离清廷，帝国行…}宣统三年（1911年），围绕「多省宣布脱离清廷，帝国行政和财政网络快速分裂」的上谕、奏折与部议在中央—地方行政链中流转。核心取证：《宣统政纪》宣统三年相关条；宫中朱批奏折、内阁大库题本与《清史稿》相关志传。这一锚点记录的是文书链：它能证明呈报、上谕、部议或照会进入机构流程，不能由此推出实际执行。来源保留：此行仅确认相关文书链和可追查的行政动作；不得由其直接推断政策有效、地方服从、民众态度或单一因果。
\eventanchor{EQ-1911-YUAN-SHIKAI-RETURN-DOCUMENTS}{围绕「清廷起用袁世凯处理北方军…}宣统三年（1911年），围绕「清廷起用袁世凯处理北方军政和议和事务」的上谕、奏折与部议在中央—地方行政链中流转。核心取证：《宣统政纪》宣统三年相关条；宫中朱批奏折、内阁大库题本与《清史稿》相关志传。这一锚点记录的是文书链：它能证明呈报、上谕、部议或照会进入机构流程，不能由此推出实际执行。来源保留：此行仅确认相关文书链和可追查的行政动作；不得由其直接推断政策有效、地方服从、民众态度或单一因果。
\eventanchor{EQ-1912-REPUBLIC-FOUNDING-DOCUMENTS}{围绕「中华民国临时政府在南京成…}宣统四年（1912年），围绕「中华民国临时政府在南京成立，与清廷退位谈判并行」的上谕、奏折与部议在中央—地方行政链中流转。核心取证：《宣统政纪》宣统四年相关条；宫中朱批奏折、内阁大库题本与《清史稿》相关志传。这一锚点记录的是文书链：它能证明呈报、上谕、部议或照会进入机构流程，不能由此推出实际执行。来源保留：此行仅确认相关文书链和可追查的行政动作；不得由其直接推断政策有效、地方服从、民众态度或单一因果。
\eventanchor{EQ-1912-QING-ABDICATION-DOCUMENTS}{围绕「宣统帝颁布退位诏书，清帝…}宣统四年（1912年），围绕「宣统帝颁布退位诏书，清帝国终结」的上谕、奏折与部议在中央—地方行政链中流转。核心取证：《宣统政纪》宣统四年相关条；宫中朱批奏折、内阁大库题本与《清史稿》相关志传。这一锚点记录的是文书链：它能证明呈报、上谕、部议或照会进入机构流程，不能由此推出实际执行。来源保留：此行仅确认相关文书链和可追查的行政动作；不得由其直接推断政策有效、地方服从、民众态度或单一因果。
\eventanchor{EQ-1910-ZIZHENG-YUAN-PRELUDE}{「资政院开院，中央层面的立宪咨…}宣统二年（1910年），「资政院开院，中央层面的立宪咨询机构开始运作」之前的条件、信息和争议已通过中央与地方材料显现，构成该事件可追踪的前史节点。核心取证：《宣统政纪》宣统二年相关条；宫中朱批奏折、内阁大库题本与《清史稿》相关志传；同时期地方档案、地方志、日记或对方材料应按具体地区补检。这一锚点界定可回查的前史条件；条件、传闻或信息累积不应被写成已经证实的单一直接原因。来源保留：官修编年和地方材料的记录密度与立场并不相同；本行不将条件、传闻、呈报或后见解释直接等同为确定因果。
\eventanchor{EQ-1910-ZIZHENG-YUAN-DECISION}{围绕「资政院开院，中央层面的立…}宣统二年（1910年），围绕「资政院开院，中央层面的立宪咨询机构开始运作」的命令、调度、照会或呈报形成独立的中央—地方行动文书链。核心取证：《宣统政纪》宣统二年相关条；宫中朱批奏折、内阁大库题本与《清史稿》相关志传。这一锚点定位决策或调度动作；命令发出、地方接收、执行过程和实际结果仍须分别寻找证据。来源保留：奏折、上谕、部议、军机档或条约可证明行政动作和机构立场，不自动证明所有地区、群体或对手按其意图行动。
\eventanchor{EQ-1910-ZIZHENG-YUAN-AFTERMATH}{「资政院开院，中央层面的立宪咨…}宣统二年（1910年），「资政院开院，中央层面的立宪咨询机构开始运作」引发的善后、执行或地方回应构成可由不同材料追踪的后续节点。核心取证：《宣统政纪》宣统二年相关条；宫中朱批奏折、内阁大库题本与《清史稿》相关志传；相关地区的档案、志书、契约、报刊或对方材料。这一锚点追踪后续影响；不同地区、群体与时段的反应不能由战后或改革后的概括性记忆代替。来源保留：战后、改革后或条约后的记忆、地方记录和官方总结常具有事后选择性；后续影响须按地点、群体和时间分层，不作整体化推断。
\eventanchor{EQ-1910-FAMINE-AND-UNREST-PRELUDE}{「灾荒、物价和财政压力持续影响…}宣统二年（1910年），「灾荒、物价和财政压力持续影响晚清地方社会与政治动员」之前的条件、信息和争议已通过中央与地方材料显现，构成该事件可追踪的前史节点。核心取证：《宣统政纪》宣统二年相关条；地方志、督抚奏折及地方档案；同时期地方档案、地方志、日记或对方材料应按具体地区补检。这一锚点界定可回查的前史条件；条件、传闻或信息累积不应被写成已经证实的单一直接原因。来源保留：官修编年和地方材料的记录密度与立场并不相同；本行不将条件、传闻、呈报或后见解释直接等同为确定因果。
\eventanchor{EQ-1910-FAMINE-AND-UNREST-DECISION}{围绕「灾荒、物价和财政压力持续…}宣统二年（1910年），围绕「灾荒、物价和财政压力持续影响晚清地方社会与政治动员」的命令、调度、照会或呈报形成独立的中央—地方行动文书链。核心取证：《宣统政纪》宣统二年相关条；地方志、督抚奏折及地方档案。这一锚点定位决策或调度动作；命令发出、地方接收、执行过程和实际结果仍须分别寻找证据。来源保留：奏折、上谕、部议、军机档或条约可证明行政动作和机构立场，不自动证明所有地区、群体或对手按其意图行动。
\eventanchor{EQ-1910-FAMINE-AND-UNREST-AFTERMATH}{「灾荒、物价和财政压力持续影响…}宣统二年（1910年），「灾荒、物价和财政压力持续影响晚清地方社会与政治动员」引发的善后、执行或地方回应构成可由不同材料追踪的后续节点。核心取证：《宣统政纪》宣统二年相关条；地方志、督抚奏折及地方档案；相关地区的档案、志书、契约、报刊或对方材料。这一锚点追踪后续影响；不同地区、群体与时段的反应不能由战后或改革后的概括性记忆代替。来源保留：战后、改革后或条约后的记忆、地方记录和官方总结常具有事后选择性；后续影响须按地点、群体和时间分层，不作整体化推断。
\eventanchor{EQ-1911-SICHUAN-RAILWAY-PROTEST-PRELUDE}{「四川保路运动升级，赵尔丰处置…}宣统三年（1911年），「四川保路运动升级，赵尔丰处置引发更大政治危机」之前的条件、信息和争议已通过中央与地方材料显现，构成该事件可追踪的前史节点。核心取证：《宣统政纪》宣统三年相关条；地方志、督抚奏折及地方档案；同时期地方档案、地方志、日记或对方材料应按具体地区补检。这一锚点界定可回查的前史条件；条件、传闻或信息累积不应被写成已经证实的单一直接原因。来源保留：官修编年和地方材料的记录密度与立场并不相同；本行不将条件、传闻、呈报或后见解释直接等同为确定因果。
\eventanchor{EQ-1911-SICHUAN-RAILWAY-PROTEST-DECISION}{围绕「四川保路运动升级，赵尔丰…}宣统三年（1911年），围绕「四川保路运动升级，赵尔丰处置引发更大政治危机」的命令、调度、照会或呈报形成独立的中央—地方行动文书链。核心取证：《宣统政纪》宣统三年相关条；地方志、督抚奏折及地方档案。这一锚点定位决策或调度动作；命令发出、地方接收、执行过程和实际结果仍须分别寻找证据。来源保留：奏折、上谕、部议、军机档或条约可证明行政动作和机构立场，不自动证明所有地区、群体或对手按其意图行动。
\eventanchor{EQ-1911-SICHUAN-RAILWAY-PROTEST-AFTERMATH}{「四川保路运动升级，赵尔丰处置…}宣统三年（1911年），「四川保路运动升级，赵尔丰处置引发更大政治危机」引发的善后、执行或地方回应构成可由不同材料追踪的后续节点。核心取证：《宣统政纪》宣统三年相关条；地方志、督抚奏折及地方档案；相关地区的档案、志书、契约、报刊或对方材料。这一锚点追踪后续影响；不同地区、群体与时段的反应不能由战后或改革后的概括性记忆代替。来源保留：战后、改革后或条约后的记忆、地方记录和官方总结常具有事后选择性；后续影响须按地点、群体和时间分层，不作整体化推断。
\eventanchor{EQ-1911-WUCHANG-UPRISING-PRELUDE}{「武昌起义爆发，新军与革命团体…}宣统三年（1911年），「武昌起义爆发，新军与革命团体控制武汉部分地区」之前的条件、信息和争议已通过中央与地方材料显现，构成该事件可追踪的前史节点。核心取证：《宣统政纪》宣统三年相关条；军机处档、战事方略及地方奏报；同时期地方档案、地方志、日记或对方材料应按具体地区补检。这一锚点界定可回查的前史条件；条件、传闻或信息累积不应被写成已经证实的单一直接原因。来源保留：官修编年和地方材料的记录密度与立场并不相同；本行不将条件、传闻、呈报或后见解释直接等同为确定因果。
\eventanchor{EQ-1911-WUCHANG-UPRISING-DECISION}{围绕「武昌起义爆发，新军与革命…}宣统三年（1911年），围绕「武昌起义爆发，新军与革命团体控制武汉部分地区」的命令、调度、照会或呈报形成独立的中央—地方行动文书链。核心取证：《宣统政纪》宣统三年相关条；军机处档、战事方略及地方奏报。这一锚点定位决策或调度动作；命令发出、地方接收、执行过程和实际结果仍须分别寻找证据。来源保留：奏折、上谕、部议、军机档或条约可证明行政动作和机构立场，不自动证明所有地区、群体或对手按其意图行动。
